\subsection{H9-T26: the full generic elliptic triple cover by Radon transform}
\label{h9:T26}

\paragraph{Status and hypotheses.}
This is a new deduction from the explicitly cited Radon and
Picard--Lefschetz theorems. Let $k$ be algebraically closed of
arbitrary characteristic, let $A/k$ be an abelian variety of
even dimension $n=2m\ge4$, and let
$f:\widetilde A=\operatorname{Bl}_B A\to\mathbf P^1_t$
be a genuine Lefschetz pencil as in H9-T19. Write $\Sigma$
for its critical values. Fix a point $P\in\mathbf P^2$ and
identify the pencil of lines through $P$ with $\mathbf P^1_t$.
Let $U$ be the open in the projective space $\mathbf P^9$ of
plane cubics consisting of smooth cubics avoiding $P$ whose
projection $\pi:C\to\mathbf P^1_t$ is unramified above
$\Sigma$. This open is nonempty and geometrically irreducible.
It represents every smooth separable degree-three genus-one
cover unramified above $\Sigma$ over algebraically closed fields,
after a choice of
plane coordinates preserving the target pencil.

\paragraph{Theorem H9-T26.}
Put $F=\overline{k(U)}$, $K=F(t)$ and
$L=F(C_{\overline\eta})$. Every homologically trivial rational
codimension-$m$ cycle on the generic fibre over $L$ is admissible,
and
\[
 h_L(z,w)=\frac13h_K(\operatorname{cor}_{L/K}z,
                             \operatorname{cor}_{L/K}w),\qquad
 \operatorname{rad}h_L=
     \operatorname{cor}_{L/K}^{-1}(\operatorname{rad}h_K).
                                                               \tag{H9.159}
\]
Thus its height quotient is the original quotient with its
form divided by three. For $m=2$, H9-T8 and accepted 032 imply
semipositivity on the full homologically trivial height domain.
The assertion is about the full geometric generic cover;
it makes no assertion for every special elliptic triple cover.

\begin{proof}
We first prove the cohomological assertion needed for the
actual Chow norm. Fix $\ell\ne\operatorname{char}k$ and work
with $\overline{\mathbf Q}_\ell$ coefficients. On
$T=\mathbf P^1-\Sigma$, write $V$ for the variable part of
$R^{n-1}f_*\overline{\mathbf Q}_\ell$ and
$\bar V=j_*V$. The already checked SGA hypotheses and
JAC-P07 give
\[
 \operatorname{rank}V=r\ge2,\quad
 H^0(\mathbf P^1,\bar V)=H^2(\mathbf P^1,\bar V)=0,
 \quad D=H^1(\mathbf P^1,\bar V),\quad
 \dim D=q-2r,\quad q=|\Sigma|.
                                                               \tag{H9.160}
\]
The representation $V$ is absolutely irreducible after every
finite-index restriction. It is self-dual up to Tate twist;
the parabolic cup pairing on $D$ is nondegenerate. This last
assertion also follows directly by identifying parabolic
cohomology with the image of compactly supported cohomology
and using its duality pairing. At every point of $\Sigma$
the local drop is one.

\emph{A simple perverse sheaf on the plane.}
Let $b:S=\operatorname{Bl}_P\mathbf P^2\to\mathbf P^2$
and $\rho:S\to\mathbf P^1$ be the induced ruling. Set
\[
                  \mathcal P=Rb_*\rho^*\bar V[2].
                                                               \tag{H9.161}
\]
The sheaf $\rho^*\bar V[2]$ is simple perverse: it is the
smooth pullback, shifted by one, of the simple perverse sheaf
$\bar V[1]$, and the fibres of $\rho$ are connected. At $P$,
proper base change gives
$i_P^*\mathcal P=D[1]$. Its costalk is concentrated in degree
$+1$ by Verdier duality. Thus the perverse support and
cosupport conditions hold at $P$, and there is neither a
punctual subobject nor a punctual quotient. Away from $P$,
$b$ is an isomorphism. Hence $\mathcal P$ is the simple
middle extension of its restriction away from $P$.
The projective bundle formula gives
\[
 R\Gamma(\mathbf P^2,\mathcal P)
  =R\Gamma(\mathbf P^1,\bar V)[2]
       \oplus R\Gamma(\mathbf P^1,\bar V)(-1).
                                                               \tag{H9.162}
\]
Its only nonzero cohomology degrees are $-1,+1$, with groups
$D,D(-1)$ respectively.

\emph{The global constant summand of the Radon transform.}
Embed $\mathbf P^2$ in $\mathbf P^a$, $a=9$, by the cubic
Veronese map $\nu$, and put $P_0=\nu_*\mathcal P$.
For the incidence projections
$p:I\to\mathbf P^a$, $q:I\to\mathbf P^{a\vee}$, use
Laumon's convention
\[
                  R(P_0)=Rq_*p^*P_0[a-1]=\mathcal R.
                                                               \tag{H9.163}
\]
Laumon Proposition~5.7(i) and (H9.162) show that
$\mathcal R$ itself is perverse. Since $P_0$ is simple and
nonconstant, Corollary~5.8(i) says that $\mathcal R$ has
exactly one nonconstant simple Jordan--H\"older constituent,
with multiplicity one; all its other constituents are constant.

On $U$, finite-index irreducibility gives
$H^0(C,\pi^*\bar V)=H^2(C,\pi^*\bar V)=0$, so
\[
 \mathcal R|_U=\mathcal W[a],\quad
 \mathcal W_C=H^1(C,\pi^*\bar V),\quad
 \mathcal W=\underline D\oplus\mathcal W_{\rm new},\quad
 \mathcal W_{\rm new}=\ker\pi_*.
                                                               \tag{H9.164}
\]
Here the first summand is the actual pullback $\pi^*D$ and
$\pi_*\pi^*=3$. We justify its extension over the entire
projective parameter space, since splitting on $U$ alone
would not suffice for the Radon argument.

The projection $p$ is a $\mathbf P^{a-1}$-bundle, and therefore
\[
 \begin{split}
 H^{-a}(\mathbf P^{a\vee},\mathcal R)
  &=H^{-1}(I,p^*P_0)\\
  &=\bigoplus_{i=0}^{a-1}
       H^{-1-2i}(\mathbf P^a,P_0)(-i)=D.
 \end{split}                                                   \tag{H9.165}
\]
By Laumon Lemma~5.6(i), the canonical adjunction map
$\underline D[a]\hookrightarrow\mathcal R$ is the maximal
constant perverse subobject. The projective-bundle
identification in (H9.165) uses pullback from $\mathbf P^a$;
restricting that adjunction map to a hyperplane section is
therefore the restriction of an ambient class. On $U$ this
is precisely the map $\pi^*:D\to H^1(C,\pi^*\bar V)$.

Verdier duality supplies a global adjoint quotient
$\mathcal R\to\underline D[a]$, after the corresponding
Tate twist is identified using the nondegenerate pairing
on $D$. Its composite with the preceding injection is
invertible: on $U$, the restriction pairing is
\[
             \langle\pi^*x,\pi^*y\rangle_C
                     =3\langle x,y\rangle_{\mathbf P^1}.
                                                               \tag{H9.166}
\]
The scalar normalization of the duality is immaterial for
invertibility. An endomorphism of a constant perverse sheaf
is determined by its restriction to $U$. Consequently the
injection splits globally, and
\[
          \mathcal R=\underline D[a]\oplus\mathcal R_{\rm new}.
                                                               \tag{H9.167}
\]
The complement has no constant subobject by maximality, and
no constant quotient by duality. Corollary~5.8(i) now makes
it simple: its three-step filtration has zero constant
bottom and zero constant top. This uses the global
adjunction map, not an assumed semisimplicity theorem.

On a genus-one cover there are $3q$ rank-one drops, so
Euler characteristic gives $\operatorname{rank}\mathcal W=3q$.
Thus
\[
              \operatorname{rank}\mathcal W_{\rm new}=2q+2r>0.
                                                               \tag{H9.168}
\]
The simple sheaf $\mathcal R_{\rm new}$ has full support,
and $\mathcal W_{\rm new}$ is an irreducible local system.

\emph{Infinite monodromy, including characteristics two and three.}
An explicit irreducible nodal cubic is
\[
 XYZ=X^3+Y^3,\qquad
 [u:v]\longmapsto[u^2v:uv^2:u^3+v^3].
                                                               \tag{H9.169}
\]
Its sole singularity is $[0:0:1]$, with tangent cone $XY$.
Projection from $[1:0:0]$ is $[Y:Z]$, and on the
normalization it is $t=u/(u^3+1)$. Its derivative is
$(1-2u^3)/(u^3+1)^2$, which is not identically zero in any
characteristic. Its degree is three. Both branches of the
node map to $t=0$. A target projective transformation puts
the finite branch set and this node value away from $\Sigma$;
it is induced by a plane coordinate change preserving the
projection centre. Thus we may take a nodal cubic $C_0$
avoiding the fixed $P$, with node over a smooth pencil value,
and with its normalization unramified over $\Sigma$.
A transverse smoothing of its ordinary node in the space
of cubics has local equation $xy=s$ up to an \'etale change
of coordinates. Properness and openness allow us to shrink
the parameter so that its punctured family lies in $U$.

Let $\widetilde C_0\to C_0$ be the normalization and put
$\mathcal F_0=\pi^*\bar V$. The normalization sequence is
\[
 0\to\mathcal F_0\to\nu_*\nu^*\mathcal F_0
                       \to i_*V_{t_0}\to0.
                                                               \tag{H9.170}
\]
The normalized projection is finite separable, so its
$H^0$ and $H^2$ with these coefficients vanish by the
finite-index irreducibility of $V$. Hence (H9.170) injects
$V_{t_0}$ into $H^1(C_0,\mathcal F_0)$, and the
nearby-cycle sequence is
\[
 0\to H^1(C_0,\mathcal F_0)
  \to H^1(C_{\bar\eta},\mathcal F_{\bar\eta})
  \to V_{t_0}(-1)\to0.
                                                               \tag{H9.171}
\]
The ordinary-node Picard--Lefschetz formula, applied to
the coefficient system which is lisse near the node,
identifies the monodromy map from the last quotient to
the gluing subspace in (H9.170) with a nonzero scalar
times the identity on $V_{t_0}$. Equivalently, the local
formula is the ordinary curve formula tensored with this
stalk; its can/variation maps are the maps in the two
displayed exact sequences. Thus its logarithm satisfies
$N^2=0$ and $\operatorname{rank}N=r$. The old summand
$\underline D$ has trivial monodromy, so the same nonzero
$N$ occurs on $\mathcal W_{\rm new}$.
This proves that its image is infinite. An irreducible
representation with infinite image has no nonzero
finite-orbit vector: the span of a finite orbit would be
a nonzero invariant subspace on which a finite-index
normal subgroup acts trivially, and irreducibility would
make the whole representation finite.

\emph{The fixed ambient term and the entire trace kernel.}
The family $\mathcal C\to U$ also has
$R^1c_*\overline{\mathbf Q}_\ell$ with no nonzero finite-orbit
vector. To see this, restrict to a genuine Lefschetz pencil
of plane cubics. Such pencils exist for the Veronese cubic
embedding in every characteristic by SGA~7 XVII; SGA~7
XVIII~6.7.2 applies to the even-dimensional ambient plane.
Its variable $H^1$ is all the elliptic $H^1$, of rank two,
and a nonzero nodal transvection makes its image infinite.
Deleting the finitely many members failing the defining
open conditions of $U$ preserves its monodromy image.
Finite-orbit vectors in the universal family would remain
finite-orbit vectors on that pencil, a contradiction.

Let $\mathcal X=\widetilde A\times_{\mathbf P^1}\mathcal C$.
It is a smooth projective family over $U$, with a finite
flat degree-three map $Q:\mathcal X\to\widetilde A\times U$.
The image of $1-Q^*Q_*/3$ on its middle cohomology is
canonically identified, by the Leray filtration, with
\[
 \bigl(H^{n-1}(A)\otimes R^1c_*\overline{\mathbf Q}_\ell\bigr)
                         \oplus\mathcal W_{\rm new}.
                                                               \tag{H9.172}
\]
Here no decomposition-theorem splitting is assumed.
Condition (A) makes all $R^bf_*$ with $b\ne n-1$ constant,
including the potentially problematic $b=n$; hence there
is no singular-fibre punctual term. The two outer Leray
terms of total degree $n$ are the trace-invariant
$H^0(R^n)$ and $H^2(R^{n-2})$. The middle term is permanent,
since on a curve it has no possible nonzero differential.
Applying the trace-kernel idempotent kills both outer terms,
so its image is canonically the middle image. Finally
$R^{n-1}f_*=H^{n-1}(A)\oplus\bar V$, and pullback/trace
on the base curves give (H9.172). Each summand displayed
has no finite-orbit vector, as just proved.

\emph{Actual cycles, heights, and full generic scope.}
Every algebraic class on the geometric generic fibre is
defined over a finite extension of the parameter field and
therefore has a finite monodromy orbit. Thus every actual
rational Chow cycle $Z\in\operatorname{CH}^m(\mathcal X_F)$
satisfies
\[
                   [Z-\tfrac13Q^*Q_*Z]=0.
                                                               \tag{H9.173}
\]
This is global homological vanishing, not Chow vanishing.
Admissibility follows from H9-T12(iii), since $\pi$ is
unramified above $\Sigma$. Pushforward and pullback of
admissible closures are admissible by the projection
formula against vertical cycles. Apply (H9.173) to such
closures, then intersect and use $Q_*Q^*=3$:
\[
 h_L(z,w)=\frac1{9}(-1)^m
          \deg(Q^*Q_*Z\cdot Q^*Q_*W)
        =\frac13h_K(\operatorname{cor}z,\operatorname{cor}w).
                                                               \tag{H9.174}
\]
Corestriction is surjective over $\mathbf Q$, because its
composite with restriction is multiplication by three.
This proves the radical formula in (H9.159).

Finally, a smooth genus-one degree-three cover has
$L=\pi^*\mathcal O(1)$ of degree three. Riemann--Roch shows
that $L$ separates every length-two subscheme and is very
ample in every characteristic. Its two defining sections
extend to a basis of $H^0(C,L)$, embedding $C$ as a plane
cubic whose projection centre misses $C$. Therefore $U$
dominates the full cover parameter problem. It is a
nonempty open of $\mathbf P^9$: the nodal smoothing above
already supplies a good smooth member. This establishes
the claimed full geometric generic scope.
\end{proof}

\paragraph{Exact sources.}
Laumon, \emph{Transformation de Fourier homog\`ene},
Bull. Soc. Math. France \textbf{131}(4) (2003),
pp.~527--551, was directly read at \S1.2 (p.~529),
Proposition~1.6 (p.~532), Lemma~5.6 (pp.~547--548),
Proposition~5.7 (p.~549), and Corollary~5.8 (p.~550).
The base hypotheses allow $S=\operatorname{Spec}k$ with
$k$ any algebraically closed field; there is no finite-field
or characteristic-zero restriction here.
\url{https://www.numdam.org/item/10.24033/bsmf.2454.pdf}.
The ordinary-node formula used in (H9.171) is SGA~7~II,
XV~3.3.5--3.3.6, printed p.~192, in relative dimension
one, tensored locally with the lisse coefficient stalk.
The earlier direct SGA checks of XVII~2.5.1 and
XVIII~6.7.2 supply the indicated cubic-pencil input.
The global constant-splitting argument, the local-coefficient
node application, and their height consequence are deductions
proved here, not quotations of an existing height theorem.
